% =====================================================================================================
\section{Introdução}\label{sec:intro}

\subsection{O problema em uma frase}
Um conjunto de aeronaves hostis se aproxima de uma região a ser defendida. O defensor dispõe de um conjunto de
\emph{interceptadores} --- pequenas aeronaves capazes de alcançar e neutralizar uma ameaça --- distribuídos em bases, e
de uma rede de radares que detecta as ameaças. A cada instante, o defensor precisa decidir \emph{quais interceptadores
lançar, contra quais ameaças e quando}, sabendo que cada interceptador só pode ser usado uma vez e que as ameaças
mais perigosas são as que se dirigem às áreas mais importantes. Este artigo trata dessa decisão, que chamamos de
\emph{alocação de interceptadores}.

\subsection{Origem histórica: o problema de atribuição arma--alvo}
A pergunta ``quais armas devem engajar quais alvos?'' é um dos problemas mais antigos da pesquisa operacional, a
disciplina que aplica modelos matemáticos a decisões de alocação de recursos. Ela foi formalizada no fim da década de
1950, no contexto da defesa contra bombardeiros e mísseis, como o problema de \conceito{atribuição arma--alvo}
(\emph{weapon--target assignment}, WTA)~\cite{manne1958}. Na forma clássica, há um conjunto de armas e um conjunto de
alvos; cada alvo tem um \emph{valor} (quanto dano causa se sobreviver) e cada par arma--alvo tem uma
\emph{probabilidade de abate} (a chance de aquela arma destruir aquele alvo). Quer-se escolher qual arma vai contra qual
alvo de modo a minimizar o valor esperado dos alvos que sobrevivem.

Um exemplo mínimo dá a intuição. Suponha dois alvos, um de valor 10 e outro de valor 1, e duas armas, cada uma com
probabilidade 0,8 de destruir qualquer alvo. Mandar uma arma contra cada alvo deixa sobreviver, em média,
$10\times0{,}2+1\times0{,}2=2{,}2$. Mandar as duas contra o alvo valioso deixa sobreviver
$10\times0{,}2\times0{,}2+1=1{,}4$, porque o alvo de valor 10 só sobrevive se as duas armas falharem. A melhor decisão
depende, portanto, dos valores e das probabilidades, e não é óbvia mesmo em instâncias pequenas.

Nos anos seguintes surgiram as primeiras heurísticas --- regras práticas que dão boas soluções sem garantia de serem as
melhores ---, como o \emph{máximo retorno marginal}, que designa uma arma de cada vez ao alvo em que ela mais reduz o
valor sobrevivente~\cite{denbroeder1959}. Na década de 1980, a defesa contra mísseis balísticos motivou a versão
\conceito{dinâmica} do problema, em que os engajamentos ocorrem em estágios e o defensor observa o resultado de um
estágio antes de decidir o próximo~\cite{hosein1990}. Na mesma época, provou-se que o WTA é \conceito{NP-completo}: não
se conhece algoritmo que o resolva exatamente em tempo que cresça só polinomialmente com o número de armas e
alvos~\cite{lloyd1986}, o que, na prática, significa que instâncias grandes só podem ser resolvidas aproximadamente.
Três revisões consolidaram a área: \citet{cai2006} sobre a versão dinâmica, \citet{roux2007} sobre sistemas de avaliação
de ameaças e atribuição de armas (TEWA, na sigla em inglês) e \citet{kline2019} sobre formulações e algoritmos exatos e
heurísticos.

\subsection{O problema hoje: enxames baratos que chegam em ondas}
Na última década, o problema mudou de escala e de natureza. Aeronaves não tripuladas baratas permitem ao atacante
lançar dezenas a centenas de vetores de uma vez, e a defesa tradicional, baseada em poucos mísseis caros, deixa de ser
economicamente viável: não faz sentido gastar um míssil de alto custo contra um drone que custa uma pequena fração dele.
Alguns casos ilustram a progressão:
\begin{itemize}
  \item Em janeiro de 2018, treze pequenos drones de construção rudimentar atacaram a base aérea russa de Khmeimim e o
        ponto naval de Tartus, na Síria; segundo o relato russo, foram neutralizados por guerra eletrônica e defesa
        antiaérea~\cite{bellingcat2018}. Foi um dos primeiros ataques amplamente noticiados com vários drones
        coordenados.
  \item Em setembro de 2019, o ataque às instalações de petróleo de Abqaiq e Khurais, na Arábia Saudita, combinou 18
        drones e cerca de sete mísseis de cruzeiro e interrompeu temporariamente cerca de metade da produção saudita, com
        a defesa praticamente sem engajar~\cite{idsa2019}. Mostrou que poucas dezenas de vetores baratos podem causar dano
        estratégico.
  \item Em abril de 2024, o ataque iraniano contra Israel empregou mais de 300 vetores --- cerca de 170 drones, mais de
        30 mísseis de cruzeiro e mais de 120 balísticos ---, e, segundo as forças israelenses, 99\,\% foram interceptados por
        uma defesa multinacional coordenada~\cite{pbs2024}. Nesse caso, a coordenação de muitos efetores contra muitas
        ameaças foi o fator decisivo.
  \item Na guerra na Ucrânia, ataques noturnos com centenas de drones do tipo Shahed levaram, a partir de 2024--2025, ao
        uso em massa de \emph{drones interceptadores}: aeronaves baratas lançadas contra as ameaças a partir das
        indicações do radar~\cite{pravda2025}.
\end{itemize}
O último caso é exatamente o cenário deste artigo: muitos interceptadores baratos e \emph{descartáveis} (cada um
neutraliza uma única ameaça e é consumido), despachados a partir de um quadro de radar comum contra muitas ameaças que
chegam em \conceito{ondas} --- grupos sucessivos, separados por alguns minutos, que tentam \conceito{saturar} a defesa,
isto é, apresentar mais ameaças simultâneas do que ela consegue tratar. Nesse regime, a qualidade da \emph{alocação}
decide quantas ameaças passam: a mesma frota de interceptadores, mal distribuída, deixa passar ameaças que uma
distribuição melhor teria detido.

\subsection{Por que a alocação é difícil}
Três fatores tornam a decisão difícil.
\begin{enumerate}
  \item \textbf{Explosão combinatória.} Com 48 ameaças e 60 interceptadores, o número de maneiras de designar
        interceptadores a ameaças é astronômico (maior que $10^{60}$), e a decisão precisa ser tomada em fração de
        segundo e refeita continuamente.
  \item \textbf{Informação imperfeita.} O defensor não vê as aeronaves; vê \conceito{pistas de radar}. Uma pista é a
        estimativa, produzida por um algoritmo de rastreamento a partir das medidas dos radares, da posição e da
        velocidade de algo que se move. Essa estimativa tem erro; não traz a identidade da aeronave; desaparece quando a
        aeronave se esconde atrás do relevo; e pode ser um \conceito{falso alarme} --- uma pista que não corresponde a
        aeronave nenhuma, criada por ruído ou por reflexos do terreno. Nos nossos cenários, a maioria das pistas pode ser
        falsa.
  \item \textbf{Ameaças heterogêneas e iscas.} Um ataque real mistura vetores muito diferentes: drones de ataque de
        longo alcance, rápidos e de rota pré-programada; munições vagantes, que orbitam o alvo antes de mergulhar; drones
        FPV, lentos, rentes ao solo e pilotados, capazes de manobrar; e \conceito{iscas}, aeronaves baratas sem carga
        explosiva, com assinatura de radar parecida, lançadas para atrair interceptadores. Parte das ameaças voa baixo
        para se esconder do radar no relevo. Quase todos os métodos de alocação estimam onde a ameaça estará supondo que
        ela mantém a velocidade atual, e nenhum tem como saber se uma pista é um drone armado ou uma isca --- e cada
        interceptador gasto numa isca é um a menos para a próxima ameaça real.
\end{enumerate}

\subsection{Lacunas da literatura}
A literatura recente oferece muitas propostas, mas três lacunas impedem saber qual usar:
\begin{enumerate}
  \item \textbf{Hipóteses otimistas.} A maior parte da literatura de WTA supõe alvos conhecidos e probabilidades de abate
        dadas; trabalhos de robótica e controle supõem consciência situacional perfeita e ameaças em linha reta ou com
        velocidade constante. Nenhuma dessas hipóteses vale contra um enxame misto, com iscas, visto por radar.
  \item \textbf{Escala pequena e cenários abstratos.} Os trabalhos brasileiros mais próximos, do Instituto Tecnológico
        de Aeronáutica~\cite{ricardo2023,feitosa2025,afonso2020,cruz2025}, modelam com cuidado a dinâmica dos defensores,
        mas usam de 3 a 6 ameaças. Em geral, os cenários da literatura são sintéticos: alvos e rotas sorteados, e radares
        posicionados sabendo de onde o ataque virá.
  \item \textbf{Comparações incomparáveis.} Os trabalhos de aprendizado por reforço~\cite{liu2023drones,palmas2025,jia2026}
        e de contra-enxame~\cite{strickland2018} usam simuladores, métricas e linhas de base próprios --- muitas vezes uma
        só linha de base. Não é possível saber, a partir da literatura, se uma política aprendida supera o método húngaro,
        ou se uma árvore de comportamento supera um leilão.
\end{enumerate}

\subsection{Contribuições e organização}
Este artigo contribui com:
\begin{itemize}
  \item uma \textbf{formulação} do problema como processo de decisão parcialmente observável, decomposta em avaliação de
        ameaças, alocação e guiagem, que explicita o que cada família de métodos supõe sobre a informação e sobre o alvo
        (Seção~\ref{sec:problema});
  \item uma \textbf{revisão} das técnicas de alocação em oito famílias, da pesquisa operacional ao aprendizado por
        reforço, explicando a ideia de cada uma e comparando os trabalhos mais próximos (Seção~\ref{sec:revisao});
  \item a \textbf{implementação de todas as técnicas revisadas} numa interface comum --- 29 métodos, cada um descrito por
        origem, funcionamento, adaptações e parâmetros ---, com a verificação de que as implementações reproduzem as
        equivalências teóricas conhecidas (Seção~\ref{sec:metodos});
  \item um \textbf{conjunto de missões realistas} e aberto: 192 missões em seis cidades brasileiras, com enxames de até
        48 vetores que misturam drones de ataque de longo alcance, iscas, munições vagantes e FPV (parâmetros de fontes
        abertas), atacando ativos reais de infraestrutura do OpenStreetMap por rotas planejadas sobre o relevo, contra
        uma defesa montada sem conhecer o ataque (Seção~\ref{sec:dataset});
  \item uma \textbf{avaliação reprodutível} num simulador aberto (Seção~\ref{sec:ambiente}), com divisão treino/teste
        por cidade, sementes pareadas e testes estatísticos, que mede tanto o dano quanto o desperdício de
        interceptadores (Seções~\ref{sec:protocolo} e~\ref{sec:resultados}).
\end{itemize}
A Seção~\ref{sec:discussao} discute os resultados, as limitações e as recomendações.
